\begin{abstract}
Thin-film materials experiments require coordinated solution preparation, liquid transport, coating, and post-processing, while reliable automation must address both robotic execution and the quality of the resulting samples. We present a hybrid robotic framework organized around heterogeneous workstation manipulators and a mobile transport platform, with thin-film fabrication as the primary application. Constraint-aware classical planning supports inter-station transport, posture-aware reinforcement learning supports labware grasping, and behavioral cloning with domain randomization and regularized few-shot real-demonstration fine-tuning supports precision liquid operations. A failure-aware hierarchical state machine sequences reusable skills using explicit completion, timeout, and recovery guards. Timestamped physical-to-simulation state synchronization provides auxiliary monitoring and trajectory comparison, rather than predictive closed-loop control. The reported individual-skill evaluation achieves an 82.7\% pooled success rate with domain randomization and few-shot adaptation; this is not an end-to-end fabrication success rate. Thin-film experiments have been completed, but quantitative comparisons of manual, fixed-trajectory, and learned execution in thickness uniformity, contact angle, roughness, and corrosion performance remain to be documented. Held-out-condition transfer tests and complete-workflow quality acceptance are specified as further validation, while serial dilution is treated as an extension of the liquid-handling skill library.
\end{abstract}
